\chapter{Classical Fourier series on an interval}
\label{module:M03.19}
\SourceRef{Zorich II, sections 18.1--18.2; Apostol, Fourier-series part of
chapter 11; Rudin, section 8 and the Fejer exercise}

\KEntry{K03.19.periodic_setup}{A fixed Fourier normalization}
\begin{definition}
Let $f:\R\to\C$ be $2\pi$-periodic and Riemann integrable on one period.
For $n\in\Z$ set
\[
 c_n(f)=\frac1{2\pi}\int_{-\pi}^{\pi}f(x)e^{-inx}\dd x,
 \qquad S_N(f)(x)=\sum_{n=-N}^Nc_n(f)e^{inx}.
\]
For real $f$, put
$a_n=\pi^{-1}\int_{-\pi}^{\pi}f(x)\cos nx\dd x$ and
$b_n=\pi^{-1}\int_{-\pi}^{\pi}f(x)\sin nx\dd x$ for $n\geq1$,
and $a_0=\pi^{-1}\int f$. Then
$c_0=a_0/2$, $c_n=(a_n-ib_n)/2$, $c_{-n}=(a_n+ib_n)/2$.
\end{definition}
These formulas follow by inserting Euler's identity into the integrals.
Finite values at a period seam may be assigned consistently with periodicity;
changing such isolated values does not change any coefficient. One-sided
limits at the seam still have to be computed from both adjacent periods.

\KEntry{K03.19.orthogonality}{Trigonometric orthogonality}
\begin{proposition}
With $e_n(x)=e^{inx}$ and
$\langle f,g\rangle=(2\pi)^{-1}\int_{-\pi}^{\pi}f\overline g$,
we have $\langle e_n,e_m\rangle=\delta_{nm}$.
The real sine/cosine orthogonality formulas follow from this identity.
\end{proposition}
\begin{proof}
If $n=m$ the normalized integral is $1$. Otherwise a primitive of
$e^{i(n-m)x}$ is $e^{i(n-m)x}/(i(n-m))$, whose endpoint values are equal
because $n-m$ is an integer. Expanding sine and cosine as
$(e_n-e_{-n})/(2i)$ and $(e_n+e_{-n})/2$ gives zero mixed integrals;
the integral mean of $\cos^2 nx$ or $\sin^2 nx$ is $1/2$ for $n\geq1$.
All products are Riemann integrable. This pairing is a seminorm pairing
on arbitrary integrable functions: zero integral norm need not mean equality
at every point. No unconstructed $L^2$ quotient is being assumed.
\end{proof}

\KEntry{K03.19.finite_projection}{Finite projection and best approximation}
\begin{theorem}
For a finite $E\subset\Z$ let $P_Ef=\sum_{n\in E}c_n(f)e_n$. For every
$p=\sum_{n\in E}d_ne_n$,
\[
 \|f-p\|_2^2=\|f-P_Ef\|_2^2+\sum_{n\in E}|c_n(f)-d_n|^2,
 \qquad \|h\|_2^2=\frac1{2\pi}\int_{-\pi}^{\pi}|h|^2.
\]
\end{theorem}
\begin{proof}
For $n\in E$, linearity and orthogonality give
$\langle f-P_Ef,e_n\rangle=c_n(f)-c_n(f)=0$.
Expand $f-p=(f-P_Ef)+(P_Ef-p)$ and expand its squared absolute value
inside the integral. Both cross terms vanish by those finite orthogonality
relations, while the second square has integral
$\sum_{n\in E}|c_n-d_n|^2$. This proves the identity and its minimizing
property without any claim that the infinite trigonometric system is complete.
\end{proof}

\KEntry{K03.19.bessel}{Bessel's inequality}
\begin{theorem}
For scalar Riemann-integrable $f$ and every finite $E\subset\Z$,
$\sum_{n\in E}|c_n(f)|^2\leq\|f\|_2^2$.
Consequently $\sum_{n\in\Z}|c_n(f)|^2$ exists as a nonnegative sum and
has that bound.
\end{theorem}
\begin{proof}
Take $p=0$ in the finite projection identity to get
$\|f\|_2^2=\|f-P_Ef\|_2^2+\sum_E|c_n|^2$.
The first term is nonnegative. Increasing finite sets then have bounded
nonnegative sums, whose supremum defines the asserted countable sum.
Integrability of $|f|^2$ follows from boundedness and continuous-composition
closure of proper Riemann integrability.
\end{proof}

\KEntry{K03.19.riemann_lebesgue}{A Riemann--Lebesgue lemma with its needed extension}
\begin{theorem}
For Riemann-integrable $f$ on $[a,b]$,
$\int_a^b f(x)e^{i\lambda x}\dd x\to0$ as real $|\lambda|\to\infty$.
In particular $c_n(f)\to0$ as $|n|\to\infty$.
The same conclusion holds for a locally Riemann-integrable function on a
bounded interval with an absolutely convergent improper integral.
\end{theorem}
\begin{proof}
First take real bounded $f$. The Darboux criterion gives a partition with
$\sum_j\operatorname{osc}_{I_j}(f)|I_j|<\eps$. Choosing any tag value in
each interval defines a step function $s$ with $\int|f-s|<\eps$; endpoint
values are irrelevant. For complex $f$ do this to real and imaginary parts,
splitting the error tolerance, to obtain the same integral error bound.
For each fixed step interval,
\[
 \left|\int_u^v e^{i\lambda x}\dd x\right|\leq2/|\lambda|,
\]
so $|\int s e^{i\lambda x}|\leq2\sum_j|s_j|/|\lambda|\to0$.
The difference is bounded by $\int|f-s|<\eps$, uniformly in $\lambda$.
First let $|\lambda|\to\infty$ and then $\eps\downarrow0$.

For an absolutely convergent improper integral choose regular endpoints
whose omitted absolute tails total less than $\eps$. The first result
applies between them, and the omitted oscillatory integrals are bounded by
those same absolute tails. This proves the extension by truncation, with
no measure-theoretic convergence theorem.
\end{proof}

\KEntry{K03.19.dirichlet_kernel}{The Dirichlet kernel and convolution identity}
\begin{construction}
Define $D_N(t)=\sum_{n=-N}^Ne^{int}$. The finite geometric identity gives
\[
 D_N(t)=\frac{\sin((N+1/2)t)}{\sin(t/2)}\quad(t\notin2\pi\Z),
 \qquad D_N(0)=2N+1.
\]
It is even and real, and its integral mean on a period is $1$.
\end{construction}
\begin{proof}[Convolution formula]
Insert the coefficient integrals into the finite sum and interchange only
that finite sum with the integral:
\[
 S_N(f)(x)=\frac1{2\pi}\int_{-\pi}^{\pi}f(y)D_N(x-y)\dd y
 =\frac1{2\pi}\int_{-\pi}^{\pi}f(x-t)D_N(t)\dd t.
\]
The second equality uses a translation followed by periodicity: a periodic
integrable function has the same integral over any interval of its period,
as follows by splitting an interval at one seam and translating the other
piece. The sine quotient follows by multiplying the finite sum by
$e^{it/2}-e^{-it/2}$ and canceling all middle terms. The excluded values
are assigned by the finite sum, so no singular integral is introduced.
Finally orthogonality of the nonconstant exponentials gives
$(2\pi)^{-1}\int D_N=1$.
\end{proof}

\KEntry{K03.19.local_convergence}{Dini's local condition}
\begin{theorem}
Let $f$ be periodic Riemann integrable, with finite limits $L_-,L_+$ at
$x$ from the two sides. If for some $0<\delta<\pi$,
\[
 \int_0^\delta
 \frac{|f(x+t)+f(x-t)-L_+-L_-|}{t}\dd t<\infty,
\]
then $S_N(f)(x)\to(L_-+L_+)/2$.
In particular the conclusion holds under piecewise $C^1$ regularity with
finite one-sided derivative bounds near the point.
\end{theorem}
\begin{proof}
By evenness and normalization of $D_N$,
\[
 S_N(f)(x)-\frac{L_-+L_+}{2}
 =\frac1{2\pi}\int_0^\pi
 \frac{f(x+t)+f(x-t)-L_+-L_-}{\sin(t/2)}
          \sin((N+1/2)t)\dd t.
\]
The quotient before the sine is locally Riemann integrable for $t>0$.
Near zero, $t/\sin(t/2)$ is bounded and tends to $2$, so the assumed
absolute improper integrability proves absolute integrability of that
quotient. Away from zero it is properly integrable. Apply the improper
extension of the preceding Riemann--Lebesgue lemma with frequency
$N+1/2$. If $f$ has one-sided Lipschitz estimates, the numerator is $O(t)$,
and its quotient by $t$ is bounded near zero. Piecewise $C^1$ functions
with finite one-sided derivative bounds have these estimates by the norm
mean-value inequality. Finite values assigned at the jump do not enter
the coefficient integrals or the limits.
\end{proof}

\KEntry{K03.19.jordan}{Dirichlet--Jordan convergence for BV functions}
\begin{theorem}
If a periodic scalar function has bounded variation on a period, then at
every $x$,
$S_N(f)(x)\to(f(x-)+f(x+))/2$.
\end{theorem}
\begin{proof}
We supply the oscillatory estimate needed when Dini's absolute condition
is unavailable. The integrals of $D_N$ over arbitrary subintervals of
$[0,\pi]$ have a common absolute bound $C$ independent of $N$.
Indeed write
\[
 \frac1{\sin(t/2)}=\frac2t+r(t),
\]
where $r$ extends continuously to $0$ by the Taylor estimate for sine.
The $r$ contribution is bounded by $\int_0^\pi|r|$. The other contribution
is $2\int_{\nu u}^{\nu v}\sin y/y\dd y$, $\nu=N+1/2$, whose absolute
value is uniformly bounded. To see the latter directly, the integral on
$[0,1]$ is bounded by $1$, and integration by parts on $[1,T]$ bounds
it by the sum of the endpoint terms and $\int_1^\infty y^{-2}\dd y$.
An interval integral is the difference of two such bounded primitives.

Consider first a real BV function $h$ on $[0,\delta]$ with $h(0)=0$ and
$h(t)\to0$ as $t\downarrow0$. Jordan decomposition writes it as a
difference of nondecreasing functions. Subtract their right limits at zero
(which have equal difference zero), to obtain
$h(t)=u(t)-v(t)$ for $t>0$, where $u,v$ are nonnegative nondecreasing,
vanish at zero and tend to zero there. Given $\eps>0$, take $\delta$ so
small that $u(\delta),v(\delta)<\eps$.
The second mean-value theorem \KRef{K03.09.second_mean_value}, reflected to a nondecreasing
weight, gives
\[
 \int_0^\delta u(t)D_N(t)\dd t
       =u(\delta)\int_\xi^\delta D_N(t)\dd t
\]
for some $\xi\in[0,\delta]$, and the same statement for $v$.
Its hypotheses hold because the monotone weights are integrable and
$D_N$ is continuous, including its removable value at zero. Thus
$|\int_0^\delta hD_N|\leq2C\eps$ uniformly in $N$.

Apply this to $f(x+t)-f(x+)$ and $f(x-t)-f(x-)$, assigning their limiting
value zero at $t=0$. They are BV and have the required limits. On
$[\delta,\pi]$ divide these bounded integrable functions by $\sin(t/2)$
and apply the ordinary Riemann--Lebesgue lemma; the resulting integrals
tend to zero. The even convolution identity, with the two constants
subtracted as in the previous proof, now bounds the limsup of the error
by a constant times $\eps$. Send $\eps$ to zero. For complex $f$ apply
the real proof to its two BV components; there is no complex order or
complex Jordan decomposition being assumed.
\end{proof}

\KEntry{K03.19.coefficient_decay}{Differentiation and Fourier coefficients}
\begin{theorem}
Suppose $f\in C^k([-\pi,\pi],\C)$, $k\geq1$, and
$f^{(j)}(-\pi)=f^{(j)}(\pi)$ for $0\leq j<k$. Then
\[
 c_n(f^{(k)})=(in)^k c_n(f),\qquad
 |c_n(f)|\leq\frac1{2\pi|n|^k}\int_{-\pi}^{\pi}|f^{(k)}(x)|\dd x
 \quad(n\ne0).
\]
In fact $c_n(f)=o(|n|^{-k})$.
\end{theorem}
\begin{proof}
For one derivative integrate $f'e^{-inx}$ by parts. The boundary term
$[f(x)e^{-inx}]_{-\pi}^{\pi}$ is zero, because the values of $f$ match
and $e^{in\pi}=e^{-in\pi}=(-1)^n$. Hence $c_n(f')=in c_n(f)$.
Iterate, using the stated lower-derivative endpoint equalities at each
step. The norm integral bound follows immediately; Riemann--Lebesgue
applied to the properly integrable $f^{(k)}$ sharpens the numerator to
$o(1)$. Endpoint matching is exactly what removes the boundary corrections.
\end{proof}

\KEntry{K03.19.c1_uniform}{Absolute convergence under first-derivative control}
\begin{theorem}
For $f\in C^1([-\pi,\pi],\C)$ with $f(-\pi)=f(\pi)$,
the Fourier series converges absolutely and uniformly to the continuous
periodic extension of $f$.
\end{theorem}
\begin{proof}
For $n\ne0$, $|c_n(f)|=|c_n(f')|/|n|$. Finite Cauchy--Schwarz gives
\[
 \sum_{1\leq|n|\leq N}|c_n(f)|
 \leq\left(\sum_{1\leq|n|\leq N}|c_n(f')|^2\right)^{1/2}
      \left(\sum_{n\ne0}\frac1{n^2}\right)^{1/2}.
\]
The first factor is bounded by $\|f'\|_2$ by Bessel and the second is
finite by the $p$-series test. Increasing partial sums therefore show
$\sum_{n\in\Z}|c_n(f)|<\infty$. The M-test gives a continuous periodic
uniform sum $g$. Termwise proper integration proves $c_n(g)=c_n(f)$
for every $n$. The coefficient uniqueness theorem proved below via Fejer
means therefore gives $g=f$. That theorem does not use the current
$C^1$ statement, so the forward reference is noncircular.

The same proof works for continuous periodic $f$ that is $C^1$ on a finite
partition with finite one-sided derivative limits: integrate by parts on
each piece, where internal boundary values of $f$ cancel. The resulting
piecewise derivative is properly Riemann integrable, and the Bessel
argument is unchanged.
\end{proof}

\KEntry{K03.19.fejer_kernel}{Fejer means and a positive kernel}
\begin{construction}
Put
\[
 \sigma_N(f)=\frac1{N+1}\sum_{k=0}^N S_k(f),\qquad
 F_N(t)=\frac1{N+1}\sum_{k=0}^ND_k(t).
\]
Then
\[
 F_N(t)=\frac1{N+1}\left|\sum_{j=0}^Ne^{ijt}\right|^2
       =\frac{\sin^2((N+1)t/2)}{(N+1)\sin^2(t/2)}
\]
away from $2\pi\Z$, with value $N+1$ there. It is nonnegative and has
integral mean $1$. Its normalized mass outside any fixed
$(-\delta,\delta)\subset(-\pi,\pi)$ tends to zero.
\end{construction}
\begin{proof}
Expand the squared finite sum. The coefficient of $e^{int}$ is
$(N+1-|n|)/(N+1)$ for $|n|\leq N$, exactly its coefficient in the average
of the $D_k$. The geometric-sum identity gives the sine quotient.
Orthogonality computes the integral mean as its constant coefficient $1$.
For $\delta\leq|t|\leq\pi$,
\[
 0\leq F_N(t)\leq\frac1{(N+1)\sin^2(\delta/2)},
\]
so even its normalized integral over that set has this upper bound.
Finite averaging of the convolution identities yields
$\sigma_N(f)(x)=(2\pi)^{-1}\int f(x-t)F_N(t)\dd t$.
\end{proof}

\KEntry{K03.19.fejer}{Uniform Fejer convergence}
\begin{theorem}
For continuous $2\pi$-periodic scalar $f$, $\sigma_N(f)\to f$ uniformly
on $\R$. In particular trigonometric polynomials are uniformly dense in
continuous periodic functions.
\end{theorem}
\begin{proof}
Such an $f$ is bounded, say $|f|\leq M$, and uniformly continuous on
$\R$: reduce one point by a period into $[-\pi,\pi]$, put the other
nearby point in $[-\pi-1,\pi+1]$, and use uniform continuity on this
compact interval. Choose $\delta$ so $|f(x-t)-f(x)|<\eps$ for
$|t|<\delta$, for all $x$. Positivity and normalization of the kernel give
\[
 |\sigma_N(f)(x)-f(x)|
 \leq\eps+\frac{2M}{(N+1)\sin^2(\delta/2)}.
\]
First let $N\to\infty$ and then $\eps\downarrow0$. Every $\sigma_N$ is a
trigonometric polynomial by its finite definition. This proves density
and makes no assertion about uniform convergence of the raw sums $S_N$.
\end{proof}

\KEntry{K03.19.uniqueness}{Uniqueness of coefficients for continuous functions}
\begin{corollary}
A continuous periodic function whose Fourier coefficients all vanish is zero.
Two continuous periodic functions with the same coefficients are equal.
\end{corollary}
\begin{proof}
All $S_N$, hence all $\sigma_N$, of the zero-coefficient function are
zero. Uniform Fejer convergence identifies their limit with the function.
Apply this to the difference of two functions for the second assertion.
The continuity assumption matters when equality is pointwise: changing
a single value of an integrable periodic function changes no coefficient.
\end{proof}

\KEntry{K03.19.parseval}{Parseval for continuous periodic functions}
\begin{theorem}
For continuous periodic $f$,
\[
 \sum_{n\in\Z}|c_n(f)|^2=\frac1{2\pi}\int_{-\pi}^{\pi}|f(x)|^2\dd x.
\]
\end{theorem}
\begin{proof}
The Fejer polynomial $\sigma_N$ lies in the span of $e_n$ for $|n|\leq N$.
Finite best approximation therefore gives
$\|f-S_N(f)\|_2^2\leq\|f-\sigma_N(f)\|_2^2$.
The right side is at most the square of its uniform norm and tends to zero.
The finite orthogonal identity gives
\[
 \|f\|_2^2=\sum_{|n|\leq N}|c_n(f)|^2+\|f-S_N(f)\|_2^2.
\]
Pass to the limit of the nonnegative coefficient sums. This proves the
formula in the stated Riemann/continuous setting, without assuming a
completed Hilbert space or Plancherel theorem.
\end{proof}

\KEntry{K03.19.examples}{Three Fourier examples and justified numerical sums}
\begin{example}
For the periodic sawtooth $f(x)=x$ on $(-\pi,\pi)$, assigned value $0$
at the seams, integration by parts gives
\[
 f(x)\sim2\sum_{n=1}^\infty\frac{(-1)^{n+1}}n\sin nx.
\]
For the square wave $g(x)=\operatorname{sgn}(\sin x)$, with value zero at
its jumps, integration over $(0,\pi)$ gives
\[
 g(x)\sim\frac4\pi\sum_{j=0}^\infty\frac{\sin((2j+1)x)}{2j+1}.
\]
Both are BV on a period. Dirichlet--Jordan proves equality with their
ordinary values at continuity points and with the midpoint value $0$ at
the jumps. At $x=\pi/2$ this gives the ordinary convergent sum
$\sum_{j\geq0}(-1)^j/(2j+1)=\pi/4$.
\end{example}
\begin{proof}[Coefficient computations and a polynomial example]
Oddness kills the cosine coefficients of both odd examples. For the
sawtooth,
$\int_0^\pi x\sin nx\dd x=-\pi(-1)^n/n$; multiplying by $2/\pi$ gives
the stated $b_n$. For the square wave the same factor multiplies
$\int_0^\pi\sin nx\dd x=(1-(-1)^n)/n$, giving zero for even $n$ and
$4/(\pi n)$ for odd $n$.

The periodic extension of $h(x)=x^2$ on $[-\pi,\pi]$ is continuous
and piecewise $C^1$. Direct integration gives $a_0/2=\pi^2/3$, $b_n=0$.
Twice integrating $x^2\cos nx$ by parts gives
$a_n=4(-1)^n/n^2$. Thus
\[
 x^2=\frac{\pi^2}3+4\sum_{n=1}^\infty\frac{(-1)^n}{n^2}\cos nx
 \quad(-\pi\leq x\leq\pi).
\]
The coefficients are absolutely summable, so the sum is uniform; its
identification with $h$ also follows from the already proved BV theorem.
At $x=\pi$ this yields $\sum_{n\geq1}n^{-2}=\pi^2/6$.
Parseval is applicable to this continuous $h$, and gives
\[
 \frac{\pi^4}{5}=\frac{\pi^4}{9}
                 +8\sum_{n=1}^\infty\frac1{n^4},
\]
hence $\sum n^{-4}=\pi^4/90$. We have not applied the
continuous-function Parseval theorem to the discontinuous square wave.
\end{proof}

\KEntry{K03.19.normalization_bridge}{Periods and real/complex representations}
\begin{proposition}
For period $T>0$, use
$\widehat f_T(n)=T^{-1}\int_{a}^{a+T}f(x)e^{-2\pi inx/T}\dd x$.
If $g(t)=f(Tt/(2\pi))$, then $c_n(g)=\widehat f_T(n)$ when both
origins are zero; translating the origin changes the expected phase.
\end{proposition}
\begin{proof}
The substitution $x=Tt/(2\pi)$ changes the normalized integral on
$[-\pi,\pi]$ into $T^{-1}\int_{-T/2}^{T/2}$ of the stated integrand.
It is $T$-periodic because $n$ is an integer, so the integral may be moved
to any interval $[a,a+T]$ without change. For $f_a(x)=f(x+a)$ a direct
substitution gives $\widehat{f_a}_T(n)=e^{2\pi ina/T}\widehat f_T(n)$.
Expanding $e^{-inx}=\cos nx-i\sin nx$ proves the real/complex
coefficient conversions stated at the start; combining the $n$ and $-n$
terms proves equality of the corresponding finite partial sums. Convergence
claims transfer only through those proved identities and their stated limits.
\end{proof}

\KEntry{K03.19.fourier_boundary}{The boundary with general Fourier analysis}
\begin{remark}[Routing record]
General Fourier transforms, Plancherel, arbitrary $L^p$ convergence,
distributions and quantitative Gibbs theory belong to D11/D24.
The present finite-interval results do not assert those theories.
In particular continuity alone is not offered as a theorem guaranteeing
pointwise convergence of raw Fourier partial sums. Fejer convergence,
absolute convergence under derivative control and the explicit Dini/BV
criteria remain distinct conclusions with distinct hypotheses.
\end{remark}
